% ----------------------------------------------------------------------
% Tool lemma: monotone sandwich for the network values pi_M.
% This file is an *artifact* for Section 4.3 of the thesis (tools for the
% repair). It records that the network values pi_M decrease in M and
% always dominate pi, so that pi_M -> pi holds if and only if for every
% epsilon > 0 some network class gets within epsilon of pi.
% When integrating into the final thesis, strip the standalone preamble
% and include this under the same numbering as the surrounding document
% (keep the "% draft: ..." comment to preserve the original label).
% Cross-references below are faked (hardcoded) so this artifact compiles
% without defining the targets; comments say what they should point to.
% ----------------------------------------------------------------------
\documentclass[11pt,a4paper]{report}
\usepackage[T1]{fontenc}
\usepackage[utf8]{inputenc}
\usepackage{amsmath,amssymb,amsthm}
\usepackage{enumitem}

\numberwithin{equation}{chapter}

\theoremstyle{plain}
\newtheorem{proposition}{Proposition}[chapter]
\newtheorem{lemma}[proposition]{Lemma}
\newtheorem{theorem}[proposition]{Theorem}
\newtheorem{corollary}[proposition]{Corollary}
\theoremstyle{definition}
\newtheorem{assumption}[proposition]{Assumption}
\newtheorem{definition}[proposition]{Definition}
\newtheorem{example}[proposition]{Example}
% Upright body, bold head (not the italic amsthm "remark" style).
\newtheorem{remark}[proposition]{Remark}

\newcommand{\E}{\mathbb{E}}
\newcommand{\PP}{\mathbb{P}}
\newcommand{\R}{\mathbb{R}}
\newcommand{\N}{\mathbb{N}}
\newcommand{\F}{\mathcal{F}}
\newcommand{\HH}{\mathcal{H}}
\newcommand{\Hu}{\mathcal{H}^{\mathrm{u}}}
\newcommand{\NNet}{\mathcal{N\!N}}
\newcommand{\as}{\ \PP\text{-a.s.}}

\begin{document}

\setcounter{chapter}{3}

\chapter{Artifact: monotone sandwich for the network values}

\subsection*{The network values decrease to a limit dominating $\pi$}

% FAKE REF: the definitions of pi and pi_M referred to in prose below
%   live in Section 4.1 of the thesis (unconstrained rewriting via
%   H o delta and the approximate problem over the network classes H_M).
Recall the hedging functional and its network approximation,
\[
  \pi(X)=\inf_{\delta\in\Hu}\rho\bigl(X+W(H\circ\delta)\bigr),
  \qquad
  \pi_M(X)=\inf_{\delta\in\HH_M}\rho\bigl(X+W(H\circ\delta)\bigr),
\]
where $\HH_M\subseteq\HH_{M+1}$ are the nested network strategy classes.
The following elementary observation reduces the approximation question
to a one-sided estimate.

% draft: Lemma 4.1 (monotone sandwich for pi_M);
% same statement as lem:monotone in extension_of_approximation_on_general.tex.
\begin{lemma}\label{lem:monotone}
  For every $X$ and every $M\in\N$, one has
  \[
    \pi_M(X)\ge\pi_{M+1}(X)\ge\pi(X).
  \]
  In particular, $\lim_{M\to\infty}\pi_M(X)$ exists in $[-\infty,\infty]$
  and dominates $\pi(X)$.
\end{lemma}

\begin{proof}
  By definition, strategies in $\HH_M$ are processes
  $\delta=(\delta_k)_{k=0,\dots,n-1}$ with
  \[
    \delta_k=F_k(I_0,\dots,I_k),\qquad k=0,\dots,n-1,
  \]
  for network maps $F_k$, hence each $\delta_k$ is $\F_k$-measurable and
  $\delta$ is adapted. With the nested exhaustion of the network classes
  one has $\HH_M\subseteq\HH_{M+1}\subseteq\Hu$, so the infima defining
  $\pi_M$, $\pi_{M+1}$ and $\pi$ are taken over increasing sets, which
  gives both inequalities. The limit exists because the sequence
  $(\pi_M(X))_{M\in\N}$ is non-increasing and bounded below by $\pi(X)$
  in $[-\infty,\infty]$.
\end{proof}

\end{document}
